\chapter{相关工作}\label{chap:related-work}

\section{模仿学习}

模仿学习旨在利用专家示范数据学习策略，使机器人在相似状态下产生接近专家的动作。经典行为克隆将专家轨迹分解为状态--动作样本，并通过监督学习拟合策略函数\citep{pomerleau1989alvinn,argall2009survey}。该方法实现直接、训练稳定，在示教覆盖充分且闭环偏移较小的任务中可以取得良好效果。然而，机器人执行过程具有闭环性质：策略输出会改变后续观测分布，一旦进入训练集中稀少的状态区域，误差可能随时间累积并导致任务失败。

针对分布偏移和时间依赖问题，已有研究通常从数据和模型两方面改进行为克隆。一方面，可以通过扰动采样、交互式数据聚合或合成专家数据扩展训练分布；另一方面，可以引入历史观测、目标条件和多步动作预测，使策略不再依赖单帧状态做出短视决策。动态目标抓取尤其需要这类时间信息，因为目标速度、相对相位和闭爪时机都无法由静态位置完全决定。本文比较的BC、BC-RNN、ACT和扩散策略均属于监督模仿学习范式，但分别代表前馈回归、循环序列建模、Transformer动作分块和生成式动作建模四种不同设计。

\section{动作分块与生成式策略}

机器人操作任务中常存在一对多映射：同一观测附近可能存在多个合理动作，例如继续跟踪、提前拦截或准备闭爪。简单回归模型在这种情况下容易预测多个模式的平均值，造成动作保守或相位错误。动作分块方法通过一次预测未来一段动作序列，将短时规划结构显式纳入策略输出，有助于缓解逐步回归造成的抖动和延迟。ACT类方法进一步使用Transformer建模观测上下文与动作序列内部依赖，使策略可以在固定窗口内表达较长时间的执行意图。

扩散模型为动作生成提供了另一种思路。其核心是先向数据逐步加入噪声，再训练神经网络学习反向去噪过程\citep{ho2020denoising,song2020denoising}。Diffusion Policy将这一思想用于机器人控制，把未来动作块作为条件生成样本，并通过重叠动作块的时间集成获得平滑闭环输出\citep{chi2023diffusion}。相较于确定性回归，扩散策略原则上可以表达多峰动作分布；相较于单步策略，动作块生成更适合需要连续对齐和时机选择的操作任务。本文采用低维目标中心化状态而非图像输入，将扩散策略与确定性动作块策略置于统一接口下比较。

\section{模型预测控制}

模型预测控制是一类滚动时域优化方法。在每个控制时刻，MPC根据系统模型预测未来状态，并求解满足约束的控制序列\citep{rawlings2017model}。在机械臂操作中，MPC能够显式考虑末端位置、姿态误差、动作平滑性、关节范围和执行速度等因素，因此常作为稳定跟踪和约束执行的核心模块。对于动态目标抓取，MPC还可以结合几何预测构造拦截或运输参考轨迹。

MPC的局限在于其性能依赖状态估计、参考生成和模型精度。当观测存在高频噪声时，相邻时刻的优化目标可能剧烈变化，导致末端参考抖动或控制量不连续；当阶段逻辑依赖阈值判断时，噪声也可能改变闭爪时机。本文采用MPC的双重角色：在真值观测下，它作为算法专家生成高质量示教；在学习策略运行时，它作为共享执行器跟踪策略输出的任务空间参考。由此，实验能够区分“动作块生成能力”和“底层约束执行能力”，并分析学习策略是否能在噪声观测下补偿直接模型控制的不足。

\section{动态目标抓取}

动态目标抓取要求机器人在目标运动过程中完成相对位姿对齐和接触决策，其难度高于静态抓取。传统方法通常依赖几何预测、轨迹优化或显式状态机，在目标状态准确且运动模式简单时具有较强可靠性。学习方法则尝试从示教数据中获得目标运动、抓取时机和控制策略之间的关系，具有适应复杂模式的潜力，但需要足够覆盖任务变化和观测扰动的数据。

本文关注平面平移与绕固定轴自旋共同作用下的动态立方体抓取。该任务同时包含连续运动跟踪、短时未来预测、阶段切换和抓取后运输，适合作为比较模型控制和模仿学习策略的实验平台。与直接端到端输出关节命令不同，本文采用目标中心化任务空间参考作为学习输出，并由MPC执行器负责低层控制。这种分层设计兼顾了学习策略对动态模式的拟合能力和模型控制对执行约束的可解释性。
